%% kapa: exam=2 round=27 year=2016 period=1 subject=appraisal_practice subject_ko=감정평가실무 seq=1 number=1 points=40
%% source: Q-Net 감정평가사 기출문제 (cst00309) artlSeq=5185401 file=extracted/1008/5185401/제27회 감정평가사 2차시험 기출문제/제27회 감정평가사 2차 시험문제지/(1교시 시험문제지)감정평가실무.pdf
%% conversion: kordoc 4.24.51 -> tools/gen2.py (auto); review=proofread-gpt-sol; flags=none
\begin{longquestion}{【문제 1】}
감정평가사 甲은 부동산투자자 乙로부터 대상부동산 투자에 관한 정보의 제공을 의뢰받고, 관련 자료를 수집ㆍ분석하여 乙에게 제공하려고 한다. 제시된 자료를 참조하여 다음 물음에 답하시오. (40점) \par
\medskip\noindent 물음 1) 2016년 7월 1일 기준 대상부동산의 시장가치를 구하되, 비교방식, 수익방식, 건물을 원가법으로 하는 공시지가기준법에 의한 가격을 각각 제시하시오. (25점) \par
\medskip\noindent 물음 2) 乙이 대상부동산을 2016년 7월 1일 매입하고 3년간 보유한 후 매각한다고 했을 때, 순현재가치(NPV)를 구하고, 乙의 투자계획에 대하여 전문가로서 제시할 의견을 기술하시오. (15점) \par
<자료 1> 대상부동산 현황 \par
1. 토지현황 \par
1) 소재지 : S시 T구 W동 500번지, 501번지 \par
2) 용도지역 : 일반상업지역 \par
3) 토지특성(2필 일단의 토지로서, 부정형의 평지이며, 소로한면에 접함) \par
- 500번지 : 대, 350 m\textsuperscript{2}, 세로장방형, 평지, 소로한면 \par
- 501번지 : 대, 450 m\textsuperscript{2}, 사다리형, 평지, 소로한면 \par
2. 건물현황 \par
1) 구조 : 철근콘크리트조 슬래브지붕 \par
2) 사용승인일자 : 2005년 5월 1일 \par
3) 세부현황 : 500번지, 501번지 양 지상에 위치함 \par
{\small\centering
\resizebox{\ifdim\width>\linewidth\linewidth\else\width\fi}{!}{%
\begin{tabular}{|c|c|c|c|c|}\hline
층 별 & 구 조 & 면적(m\textsuperscript{2}) & 용 도 & 비 고 \\\hline
지하1층 & 철근콘크리트조 & 260 & 주차장, 기계실 & - \\\hline
지상1층 & 철근콘크리트조 & 520 & 사무실 & P 은행 임차 \\\hline
지상2층 & 철근콘크리트조 & 520 & 사무실 & P 은행 임차 \\\hline
지상3층 & 철근콘크리트조 & 520 & 사무실 & R 회사 임차 \\\hline
지상4층 & 철근콘크리트조 & 520 & 사무실 & R 회사 임차 \\\hline
지상5층 & 철근콘크리트조 & 400 & 사무실 & 공 실 \\\hline
\multicolumn{2}{|c|}{계} & 2,740 & \multicolumn{2}{c|}{-} \\\hline
\end{tabular}}\par}

3. 기타현황 \par
: 현재 대상부동산의 소유자는 丙과 丁이 공동소유(각각 50 \%)를 하고 있으며, 소유자 丙의 명의로 C은행에 근저당권 5억원이 설정되어 있음 \par
<자료 2> 표준지 공시지가 현황(공시기준일 : 2016.01.01) \par
{\small\centering
\resizebox{\ifdim\width>\linewidth\linewidth\else\width\fi}{!}{%
\begin{tabular}{|c|c|c|c|c|c|c|c|c|c|}\hline
\makecell{기호} & \makecell{소재지\\(지번)} & \makecell{면적\\(m\textsuperscript{2})} & 지목 & \makecell{이용\\상황} & \makecell{용도\\지역} & \makecell{주위\\환경} & \makecell{도로\\교통} & \makecell{형상\\지세} & \makecell{공시지가\\(원/m\textsuperscript{2})} \\\hline
1 & \makecell{V동\\130} & \makecell{(일단지)\\535} & 대 & 주상용 & \makecell{일반\\상업} & \makecell{후면\\상가지대} & \makecell{소로\\한면} & \makecell{사다리형\\평지} & 2,700,000 \\\hline
2 & \makecell{V동\\150} & 800 & 대 & 상업용 & \makecell{일반\\상업} & \makecell{노선\\상가지대} & \makecell{소로\\한면} & \makecell{가장형\\평지} & 3,400,000 \\\hline
3 & \makecell{W동\\485} & 420 & 대 & \makecell{상업기타\\주차건물} & \makecell{일반\\상업} & \makecell{후면\\상가지대} & \makecell{세로\\(가)} & \makecell{사다리형\\평지} & 2,100,000 \\\hline
4 & \makecell{W동\\520} & 450 & 대 & 상업용 & \makecell{일반\\상업} & \makecell{노선\\상가지대} & \makecell{소로\\한면} & \makecell{사다리형\\평지} & 3,000,000 \\\hline
\end{tabular}}\par}
<자료 3> 평가선례(공시지가기준법 적용시 그 밖의 요인에 적용함) \par
{\small\centering
\resizebox{\ifdim\width>\linewidth\linewidth\else\width\fi}{!}{%
\begin{tabular}{|c|c|c|c|}
\hline
\makecell{구 분} & \makecell{평가선례 \#1} & \makecell{평가선례 \#2} & \makecell{평가선례 \#3} \\
\hline
\makecell{소재지(지번)} & \makecell{V동 143} & \makecell{W동 504} & \makecell{W동 522} \\
\hline
\makecell{토지현황} & \makecell{대, 450 m\textsuperscript{2},\\소로한면, 부정형,\\평지, 일반상업지역} & \makecell{대, 780 m\textsuperscript{2},\\소로한면, 사다리형,\\평지, 일반상업지역} & \makecell{대, 350 m\textsuperscript{2},\\세로(가), 가장형,\\평지, 일반상업지역} \\
\hline
\makecell{건물현황} & \makecell{철근콘크리트조\\슬래브지붕, 상업용,\\지하1층∼지상5층,\\연면적 3,300 m\textsuperscript{2}} & \makecell{철근콘크리트조\\슬래브지붕, 상업용,\\지하2층∼지상8층,\\연면적 5,200 m\textsuperscript{2}} & \makecell{없음(상업나지)} \\
\hline
\makecell{토지단가} & \makecell{3,600,000원/m\textsuperscript{2}} & \makecell{3,400,000원/m\textsuperscript{2}} & \makecell{3,500,000원/m\textsuperscript{2}} \\
\hline
\makecell{기타사항} & \makecell{기준시점 2016.04.01\\일반거래목적의\\정상적인 평가선례} & \makecell{기준시점 2016.01.01\\일반거래목적의\\정상적인 평가선례} & \makecell{기준시점 2016.05.01\\담보목적의\\정상적인 평가선례} \\
\hline
\end{tabular}}\par}

<자료 4> 거래사례(거래사례비교법에 적용함) \par
{\small\centering
\resizebox{\ifdim\width>\linewidth\linewidth\else\width\fi}{!}{%
\begin{tabular}{|c|c|c|c|}
\hline
\makecell{구 분} & \makecell{사례 \#1} & \makecell{사례 \#2} & \makecell{사례 \#3} \\
\hline
\makecell{소재지(지번)} & \makecell{V동 135번지} & \makecell{W동 489번지} & \makecell{W동 515번지} \\
\hline
\makecell{토지현황} & \makecell{대, 900 m\textsuperscript{2},\\소로한면, 부정형,\\평지, 일반상업지역} & \makecell{대, 550 m\textsuperscript{2},\\세로가, 사다리형,\\평지, 일반상업지역} & \makecell{대, 750 m\textsuperscript{2},\\소로한면, 가장형,\\평지, 일반상업지역} \\
\hline
\makecell{건물현황} & \makecell{철근콘크리트조\\슬래브지붕, 상업용,\\지하1층∼지상5층,\\연면적 3,200 m\textsuperscript{2},\\2006.05.01 사용승인} & \makecell{없음(상업나지)} & \makecell{철근콘크리트조\\슬래브지붕, 상업용,\\지상2층,\\연면적 800 m\textsuperscript{2},\\1985.01.01 사용승인} \\
\hline
\makecell{거래금액} & \makecell{5,600,000,000원} & \makecell{1,600,000,000원} & \makecell{2,850,000,000원} \\
\hline
\makecell{기타사항} & \makecell{노선상가지대\\2016.01.01 거래\\정상거래사례\\(토지건물가격비중 6:4)} & \makecell{후면상가지대\\2016.04.01 거래\\정상적인 거래사례} & \makecell{노선상가지대\\2016.05.01 거래\\철거전제의 정상거래,\\매수자의 철거비 부담\\(30,000,000원)} \\
\hline
\end{tabular}}\par}

<자료 5> 표준건축비 등 \par
1. 인근지역 상업용건축물(철근콘크리트조 5층 이하)의 표준건축비(부대설비 포함, 지상ㆍ지하 건축물에 동일하게 적용)는 770,000원/m\textsuperscript{2}이며, 건물의 잔존가치는 10 \%임 \par
<자료 6> 대상부동산의 임대내역 \par
{\small\centering
\resizebox{\ifdim\width>\linewidth\linewidth\else\width\fi}{!}{%
\begin{tabular}{|c|c|c|c|c|}
\hline
\makecell{구분} & \makecell{면적\\(m\textsuperscript{2})} & \makecell{임차인} & \makecell{임대기간} & \makecell{임대료} \\
\hline
\makecell{지상\\1층} & \makecell{520} & \makecell{P 은행} & \makecell{2011.07.01 ～ 2016.06.30} & \makecell{‘15년 7월 1일부터 연간가능총\\소득(PGI)120,000원/m\textsuperscript{2} 적용} \\
\hline
\makecell{지상\\2층} & \makecell{520} & \makecell{P 은행} & \makecell{2011.07.01 ～ 2016.06.30} & \makecell{‘15년 7월 1일부터 연간가능총\\소득(PGI)95,000원/m\textsuperscript{2} 적용} \\
\hline
\makecell{지상\\3층} & \makecell{520} & \makecell{R 회사} & \makecell{2014.07.01 ～ 2019.06.30} & \makecell{‘15년 7월 1일부터 연간가능총\\소득(PGI)80,000원/m\textsuperscript{2} 적용} \\
\hline
\makecell{지상\\4층} & \makecell{520} & \makecell{R 회사} & \makecell{2014.07.01 ～ 2019.06.30} & \makecell{‘15년 7월 1일부터 연간가능총\\소득(PGI)80,000원/m\textsuperscript{2} 적용} \\
\hline
\makecell{지상\\5층} & \makecell{400} & \makecell{공 실} & \makecell{-} & \makecell{최근 1개월간 공실} \\
\hline
\makecell{계} & \makecell{2,480} & \multicolumn{3}{|c|}{\makecell{-}} \\
\hline
\end{tabular}}\par}

※ R회사는 회사 사정상 2016.06.30에 이전할 계획이며, 현재 소유자도 중도 계약 해지에 동의하였고, 새로운 임차인을 시장임대료로 즉시 구할 수 있음 \par
<자료 7> 최근 임대사례 \par
1. 사례물건 : V동 138번지 소재 5층 \par
1) 토지현황 : 일반상업지역, 대, 950 m\textsuperscript{2}, 소로한면, 사다리형, 평지 \par
2) 건물현황 : 철근콘크리트조, 지하1층∼지상5층, 연면적 3,200 m\textsuperscript{2}, 상업용 \par
2. 임대상황 \par
1) 1～2층 임대사례 : G 은행 2016.07.01부터 5년 계약 \par
연간가능총소득(PGI) 1층 160,000원/m\textsuperscript{2}, 2층 120,000원/m\textsuperscript{2} \par
2) 3～5층 임대사례 : H 회사 2016.07.01부터 5년 계약 \par
연간가능총소득(PGI) 3～4층 100,000원/m\textsuperscript{2}, 5층 90,000원/m\textsuperscript{2} \par
<자료 8> 시점수정 등 관련자료(T구) \par
1. 지가변동률 \par
{\small\centering
\resizebox{\ifdim\width>\linewidth\linewidth\else\width\fi}{!}{%
\begin{tabular}{|c|c|}
\hline
\makecell{구 분} & \makecell{일반상업지역} \\
\hline
\makecell{2016. 01. 01 ～ 2016. 05. 31(누계)} & \makecell{1.687} \\
\hline
\makecell{2016. 04. 01 ～ 2016. 05. 31(누계)} & \makecell{0.654} \\
\hline
\makecell{2016. 05. 01 ～ 2016. 05. 31(당월)} & \makecell{0.323} \\
\hline
\end{tabular}}\par}

※ 2016년 6월부터 지가변동률 미고시, 2016년 6월 지가변동률은 직전월 자료를 적용하고, 변동률의 계산은 백분율을 기준으로 소수점 넷째자리에서 반올림함 \par
2. 임대동향조사 중 소형부동산의 자본수익률 \par
{\small\centering
\resizebox{\ifdim\width>\linewidth\linewidth\else\width\fi}{!}{%
\begin{tabular}{|c|c|}
\hline
\makecell{구 분} & \makecell{T 구} \\
\hline
\makecell{2016. 01. 01 ～ 2016. 05. 31(누계)} & \makecell{2.113} \\
\hline
\makecell{2016. 04. 01 ～ 2016. 05. 31(누계)} & \makecell{0.895} \\
\hline
\makecell{2016. 05. 01 ～ 2016. 05. 31(당월)} & \makecell{0.356} \\
\hline
\end{tabular}}\par}

3. 건축비지수 : 2015년 1월 1일 이후 보합세임 \par
<자료 9> 지역요인, 개별요인 등 품등 비교자료 \par
1. 지역요인 자료 \par
: V동과 W동은 S시 T구에 속하며, 간선도로(소로한면)를 두고 맞은편에 위치하고 있음. 최근 V동 남측에 종합유통센터가 개장함에 따라 V동 상권으로 유동인구가 증가하여 V동이 W동에 비해 지역요인이 3 \% 우세를 보이고 있음 \par
2. 개별요인 자료 \par
1) 도로접면 격차율 \par
{\small\centering
\resizebox{\ifdim\width>\linewidth\linewidth\else\width\fi}{!}{%
\begin{tabular}{|c|c|c|c|c|c|c|}
\hline
\makecell{구 분} & \makecell{광대세각} & \makecell{광대한면} & \makecell{소로한면} & \makecell{세로(가)} & \makecell{세로(불)} & \makecell{맹지} \\
\hline
\makecell{광대세각} & \makecell{1.00} & \makecell{0.95} & \makecell{0.86} & \makecell{0.81} & \makecell{0.75} & \makecell{0.72} \\
\hline
\makecell{광대한면} & \makecell{1.05} & \makecell{1.00} & \makecell{0.91} & \makecell{0.86} & \makecell{0.78} & \makecell{0.75} \\
\hline
\makecell{소로한면} & \makecell{1.16} & \makecell{1.10} & \makecell{1.00} & \makecell{0.91} & \makecell{0.86} & \makecell{0.83} \\
\hline
\makecell{세로(가)} & \makecell{1.24} & \makecell{1.18} & \makecell{1.07} & \makecell{1.00} & \makecell{0.92} & \makecell{0.89} \\
\hline
\makecell{세로(불)} & \makecell{1.34} & \makecell{1.28} & \makecell{1.16} & \makecell{1.07} & \makecell{1.00} & \makecell{0.96} \\
\hline
\makecell{맹지} & \makecell{1.39} & \makecell{1.32} & \makecell{1.20} & \makecell{1.16} & \makecell{1.04} & \makecell{1.00} \\
\hline
\end{tabular}}\par}

2) 형상 격차율 \par
{\small\centering
\resizebox{\ifdim\width>\linewidth\linewidth\else\width\fi}{!}{%
\begin{tabular}{|c|c|c|c|c|}
\hline
\makecell{구 분} & \makecell{정방형} & \makecell{장방형} & \makecell{사다리형} & \makecell{부정형} \\
\hline
\makecell{정방형} & \makecell{1.00} & \makecell{0.99} & \makecell{0.98} & \makecell{0.95} \\
\hline
\makecell{장방형} & \makecell{1.01} & \makecell{1.00} & \makecell{0.99} & \makecell{0.96} \\
\hline
\makecell{사다리형} & \makecell{1.02} & \makecell{1.01} & \makecell{1.00} & \makecell{0.97} \\
\hline
\makecell{부정형} & \makecell{1.05} & \makecell{1.04} & \makecell{1.03} & \makecell{1.00} \\
\hline
\end{tabular}}\par}

※ 가로장방형은 장방형을 적용 \par
3) 환경조건 격차율 \par
ㆍ대상부동산은 표준지 기호1 보다 5 \% 우세하며, 표준지 기호3 보다 15 \% 우세함 \par
ㆍ대상부동산은 평가선례 \#1 보다 10 \% 열세하며, 거래사례 \#2 보다 15 \% 우세함 \par
3. 임대사례(V동 138번지)와 대상부동산의 품등 격차율(지역 요인과 개별요인 포함) : 임대사례가 대상부동산 보다 총 10 \% 우세함 \par
4. 거래사례 \#1(V동 135번지)의 건축물과 대상부동산의 건축물은 동일한 등급 수준으로 신축되었음 \par
<자료 10> 수익환원법 적용 자료 및 의뢰인 乙의 부동산 투자계획 \par
1. 환원율 및 할인율 \par
1) 현재시점 환원율 : 시장추출법에 의한 산정 \par
{\small\centering
\resizebox{\ifdim\width>\linewidth\linewidth\else\width\fi}{!}{%
\begin{tabular}{|c|c|c|c|}
\hline
\makecell{구 분} & \makecell{사례 \#1} & \makecell{사례 \#2} & \makecell{사례 \#3} \\
\hline
\makecell{매매가격(원)} & \makecell{3,500,000,000} & \makecell{2,200,000,000} & \makecell{2,400,000,000} \\
\hline
\makecell{순수익(원)} & \makecell{140,000,000} & \makecell{88,000,000} & \makecell{200,000,000} \\
\hline
\makecell{기타} & \makecell{최근사례, 정상거래} & \makecell{최근사례, 정상거래} & \makecell{최근사례, 사정개입} \\
\hline
\end{tabular}}\par}

2) 재매도가치 산정을 위한 환원율 \par
: 현재 환원율에 장기위험프리미엄 등을 고려하여 0.5 \% 가산함 \par
3) 할인현금흐름분석법에 사용할 할인율 : 투자자의 요구수익률 \par
2. 수익환원법 적용 \par
1) 대상부동산의 시장가치산정을 위한 수익환원법은 1년차 순영업소득(NOI)을 직접 환원하는 직접환원법을 적용 \par
2) 수익환원법에 적용되는 수익과 비용은 연간 단위로 산정하고 연말에 인식하는 것을 가정하며, 연간가능총소득(PGI)에는 관리비 등 제반 내역이 합리적으로 포함되어 있다고 전제함 \par
3) 대상부동산의 수익환원법 적용시, 연간가능총소득(PGI)은 최근 임대사례에서 산출하고, 공실손실상당액, 운영경비 등은 대상부동산을 기준으로 산정함 \par
4) 연간가능총소득(PGI)은 매년 5 \% 상승하고, 공실률은 매년 5 \%로 예상되며, 운영경비는 각 층별 면적 기준으로 25,000원/m\textsuperscript{2}이 소요되며, 운영경비는 매년 4 \% 상승을 적용함 \par
3. 乙의 투자계획 \par
1) 투자금액 : 4,200,000,000원 \par
2) 요구수익율 : 6 \% \par
3) 투자기간 : 3년간 보유한 후 매각 \par
4) 보유기간말 매각시 매각비용 : 매각금액의 3 \% \par
4. 순현재가치(NPV)의 산정방법 \par
: 순현재가치(NPV)는 乙이 투자계획대로 대상부동산을 3년간 보유한 후 매각하는 것을 가정하여 산정하며, 수익환원법 적용 자료와 같이 연간단위로 수익과 비용을 연말에 인식함 \par
\end{longquestion}
